\subsection{Retrospective Protocol Comparison}
\label{sec:real}

Using complete expert annotations, we reconstruct the measurements that each
partial protocol would have produced and evaluate fitted predictors on
held-out units. Sleep protocols sample binary stage indicators from scored
hypnograms, whereas AF protocols average binary rhythm annotations within
selected time bins. Each reconstructed protocol is therefore defined by the
subset or aggregation of annotation times made available to the predictor.

Sleep contains \numSleepRecordings{} recordings from \numSleepSubjects{}
subjects and uses $p=\numSleepGrid$ relative-time anchors. Its targets are the
exact full-interval REM, N3 and Wake proportions, and its budgets
$N\in\{4,8,16,32,64\}$ count distinct scored 30-second epochs. Long-Term AF
contains \numLtafRecords{} records, uses $p=\numAfGrid$ equal relative-time
bins and targets the exact fraction of analysable annotated time spent in AF.
Its budgets $N\in\{1,2,4,8,16,32\}$ observe the fraction $N/p$ of each record.
Both analyses use five outer folds and report pooled held-out cross-fitted
$R^2$. Sleep folds are subject-disjoint and stratified by source study; AF
folds are split by record because repeated-subject identifiers are unavailable.

Two fixed templates are compared at every budget: a centred-contiguous block
and anchors dispersed as uniformly as possible over relative time. Sleep also
learns target-aware and kernel-quadrature supports entirely within each outer
training fold; AF uses only the pre-specified templates. The learned Sleep
procedure uses a regularised best-linear score for foldwise protocol selection
and is evaluated by held-out cross-fitted $R^2$; \cref{sec:estimation}
analyses the corresponding fixed-dimensional Gaussian plug-in setting.
Residual covariance, target cross-covariance and target variance are estimated
foldwise from the annotations, with an eigenvalue floor adapted to the
high-dimensional grid.
The common-grid mapping, baselines, weighting, selection score, nested ridge
fitting and resampling procedures are specified in
Appendix~\ref{sec:appendix-real}.

Sleep-EDF Expanded \citep{kemp2000analysis,mourtazaev1995age,
goldberger2000physiobank,pollard2026physionet} combines Sleep Cassette (SC) and
Sleep Telemetry (ST) 30-second hypnograms. SC contributes
\numSleepScRecords{} recordings from \numSleepScSubjects{} participants and ST
contributes \numSleepStRecords{} recordings from \numSleepStSubjects{}
participants. Stages 3 and 4 are combined as N3
\citep{rechtschaffen1968manual,iber2007aasm}. The retained records contain
\numSleepHours{} valid annotated hours.

For REM, \cref{fig:real}(a) shows that the centred-contiguous template has
cross-fitted $R^2$ values \numCfContigFour, \numCfContigSixteen{} and
\numCfContigSixtyFour{} at $N=4,16,64$, compared with
\numCfUniformFour, \numCfUniformSixteen{} and \numCfUniformSixtyFour{} for
uniform dispersion. Across the 15 target--budget cells, the conditional
dispersed-minus-contiguous range has a positive lower endpoint in
\numSleepAdjustedPositiveRanges{} cells.

The pooled analysis therefore favours dispersion at several budgets, but the
source-specific estimates are heterogeneous. When SC and ST are analysed by
separate subject-grouped pipelines, \numSleepCohortPositiveRanges{} of
\numSleepCohortCells{} source--target--budget ranges lie above zero,
\numSleepCohortNegativeRanges{} lie below zero and
\numSleepCohortUnresolvedRanges{} include zero. These source-specific and
fixed-template sensitivity procedures are detailed in
Appendix~\ref{sec:appendix-real-sensitivity}.


The Long-Term AF Database
\citep{petrutiu2007abrupt,goldberger2000physiobank} contains
\numLtafHours{} hours of reviewed rhythm annotations and has median analysable
record length \numLtafMedianHours{} hours. The analysable interval begins at
the first rhythm marker, so any unannotated prefix is excluded from the target
and the candidate windows. All \numAfCohortNAll{} records form the primary
cohort.

At $N=4$, each template observes \numAfObservedPctFour{} of the analysable
record, equivalent to \numAfEquivalentHoursFour{} hour on a 24-hour record.
The contiguous and dispersed templates attain cross-fitted $R^2$ values
\numCfAfContigOneHour{} and \numCfAfDispOneHour, a difference of
\numAfBootDiffFour{} with conditional paired range \numAfBootRangeFour. At
$N=16$, the fraction is \numAfObservedPctSixteen{} and the corresponding
values are \numCfAfContigFourHour{} and \numCfAfDispFourHour; their difference
is \numAfBootDiffSixteen{} with range \numAfBootRangeSixteen.

The conditional paired range has a positive lower endpoint at every evaluated
budget containing at least two windows. At $N=1$, the difference is
\numAfBootDiffOne{} with range \numAfBootRangeOne. The two single-window
templates occupy different locations, so this cell estimates a location
contrast; budgets of at least two additionally compare dispersed with
contiguous layouts. Outcome-defined cohort and grid-discretisation sensitivities
are reported in Appendix~\ref{sec:appendix-real-sensitivity}.


Every point in \cref{fig:real} is obtained by pooling the predictions from all
five outer test folds and then applying the weighted $R^2$ in
Equation~\eqref{eq:weighted-cf-r2} in
Appendix~\ref{sec:appendix-real-prediction}; foldwise $R^2$ values are not
averaged. Sleep pooling uses subject-balanced weights, whereas AF pooling gives
equal weight to each record. Within each test fold, all supports are fixed
before prediction evaluation. In panel (a), study, valid-duration and treatment
covariates enter every Sleep predictor unpenalised. In panel (b), the horizontal
coordinate is the protocol's observed fraction $N/p$ of each analysable AF
record.

\FloatBarrier
\begin{figure}[htbp]
\centering
\includegraphics[width=0.86\linewidth]{figures/fig_real.pdf}
\caption{Held-out prediction from reconstructed Sleep and AF annotation
protocols.}
\label{fig:real}
\end{figure}
\FloatBarrier

The Sleep data also allow us to examine whether the exact support locations
selected from the training data are stable.
Target-aware selection maximises the training-fold score defined in
Equation~\eqref{eq:real-objective} in
Appendix~\ref{sec:appendix-real-prediction}; its target-agnostic comparator learns a
kernel-quadrature support from the same regularised covariance without using
the target outcomes.

At $N=16$, the target-aware REM support reaches baseline-adjusted cross-fitted
$R^2=\numCfAwareSixteen$, compared with \numCfKqSixteen{} for learned kernel
quadrature and \numCfUniformSixteen{} for fixed dispersion. The fold-averaged
pairwise Jaccard overlap among REM-, N3- and Wake-specific supports is
\numSleepJaccardMin--\numSleepJaccardMax, indicating that the selected
locations depend strongly on the target. We therefore examine their stability
across support-selection sample sizes and subject subsamples.

Panel (a) of \cref{fig:sweep} varies the subjects used for support selection
from \numSweepMSmall{} to the full outer-training-fold count while leaving
ridge fitting on the complete training fold. The target-aware difference
changes from \numSweepDeltaKqSmall{} to \numSweepDeltaKqFull{} against learned
kernel quadrature and from \numSweepDeltaUniSmall{} to
\numSweepDeltaUniFull{} against fixed dispersion. Panel (b) repeats support
selection, tuning and fitting on \numSelectionSubsampleReps{} study-stratified
subsamples of \numSelectionSubsamplePct{} of subjects, drawn without
replacement. The target-aware differences have medians and 2.5--97.5 percentile
ranges \numSelectionDeltaKqMed{}
$[\numSelectionDeltaKqPlo,\numSelectionDeltaKqPhi]$ against kernel quadrature
and \numSelectionDeltaUniMed{}
$[\numSelectionDeltaUniPlo,\numSelectionDeltaUniPhi]$ against fixed
dispersion. Both ranges span zero.
Each ordinate in panel (a) is the paired difference between the target-aware
and comparator pooled cross-fitted $R^2$ for the same support-selection sample;
the plotted point averages the repeated selections and its band is one standard
error. Each panel-(b) subsample recomputes standardisation, foldwise moments,
support selection, inner ridge tuning, predictor fitting and outer-fold
evaluation before the paired differences are formed.

\begin{figure}[htbp]
\centering
\includegraphics[width=0.84\linewidth]{figures/fig_sweep.pdf}
\caption{Learned REM support at $N=16$. Each panel-(b) subsample reruns the
complete selection, tuning, fitting and held-out-evaluation pipeline.}
\label{fig:sweep}
\end{figure}
\FloatBarrier

Taken together, the pooled Sleep and AF analyses provide clearer evidence for
differences between pre-specified dispersed and contiguous layouts than for
differences among learned supports, although the Sleep
fixed-template effects remain heterogeneous across source studies. The
learned supports vary across targets and resamples, and their paired held-out
$R^2$ differences span zero. The subsampling scheme and additional weighting
and regularisation diagnostics are given in
Appendix~\ref{sec:appendix-real-sensitivity}.

